\section{Conclusion and Limitations}
\label{sec:conclusion}

We release TwinRouterBench, a step-level routing benchmark built on the
scenarios where routing matters most: long-context, multi-turn agent
workloads such as SWE-bench, PinchBench, BFCL, mtRAG, and QMSum. The
static track approximates per-step cheapest-sufficient tiers via a greedy
sequential-locking downgrade search, and a manual audit on multi-turn
trajectories confirms the reliability of the resulting labels. On the
dynamic track, a 100-case held-out SWE-bench evaluation demonstrates that
the live execution framework produces meaningful end-to-end results. A
logistic router trained on the static labels achieves a 53\% cost reduction
at matched resolve rate on the dynamic track, showing that the static data
serves not only as a fast offline evaluation tool but also as effective
training supervision for practical routers.

\paragraph{Limitations.}\label{sec:limitations}
The static corpus of 970 routed steps from 520 instances is sufficient for
initial diagnostics and training but not exhaustive across agent workloads,
domains, or routing patterns. Labels are estimated under the fixed model pool,
cascade protocol, and price frontier of \S\ref{sec:static-scoring}, and
would need updating if a future model became both cheap and highly capable.
Dynamic coverage beyond SWE-bench Verified, and keeping data current as
pools and pricing evolve, are future work.

\paragraph{Broader impacts.}
TwinRouterBench may accelerate the development of cost-effective LLM
routers by providing standardized step-level evaluation, indirectly
reducing API cost and energy consumption in agentic deployments. The
two-track design also improves reproducibility by separating
deterministic offline scoring from live validation. A potential negative
impact is benchmark overfitting: community efforts may concentrate on
the current 970-row corpus and five workloads at the expense of
uncovered domains, languages, or agent architectures. Additionally,
tier labels are tied to a fixed model pool and price snapshot; as
models and pricing evolve, stale labels could mislead router training
or evaluation. We mitigate these risks through versioned model pools
and pricing, explicit scope limitations, and a release design that
supports pool updates and re-labeling.
